\documentclass{article}
\usepackage[UTF8]{ctex}
\usepackage{amsmath,amssymb,amsfonts}
\usepackage{enumitem}
\usepackage{geometry}
\geometry{a4paper, margin=1in}
% === 错题本排版优化（不参与题目标识） ===
\setlength{\parindent}{0pt}
\setlength{\parskip}{0.35em plus 0.10em minus 0.08em}
\setlength{\emergencystretch}{3em}
\setlist{itemsep=0.25em, topsep=0.40em, parsep=0pt}
\allowdisplaybreaks[2]
% === 错题本排版优化结束 ===
\begin{document}
\section*{上海交通大学2025年数学分析真题}
\section*{43．上海交通大学2025年数学分析试卷（回忆版）}
\subsection{题目1}
1． \(\lim _{n \rightarrow \infty} \frac{1^{5}+2^{5}+\cdots+n^{5}}{n^{6}}\) ．
\subsection{题目2}
2．\(\sum_{n=1}^{\infty} \frac{n}{(n+1)!}\) ．
\subsection{题目3}
3．设 \(\Omega\) 为 \(x^{2}+y^{2}=1\) 的圆柱面在 \(0 \leq z \leq 1\) 的部分，求 \(\iint_{\Omega} \frac{\mathrm{d} S}{\sqrt{x^{2}+y^{2}+z^{2}}}\) ．
\subsection{题目4}
4．求 \(I(x)=\int_{0}^{+\infty} e^{-t^{2}} \cos (2 x t) \mathrm{d} t, x \in(-\infty,+\infty)\) ．
\subsection{题目5}
二．求证：\(\sum_{n=1}^{\infty}(-1)^{n} \frac{x^{2}+n}{n^{2}}\) 在 \(x \in[a, b]\) 上一致收敛，但对任意的 \(x\) 都不绝对收敛。
\subsection{题目6}
三．设 \(T \in(0,+\infty)\) ，存在 \(b, \tau>0\) ，对任意的 \(t \in(0, T)\) ，满足 \(\frac{1}{\tau} \int_{t}^{t+T} g(s) \mathrm{d} s<b\) ，其中 \(g(s)\) 为非负连续函数，求证：对任意的正数 \(a\) ，有 \(\int_{0}^{t} e^{-a(t-s)} g(s) \mathrm{d} s \leq \frac{b \tau}{1-e^{-a \tau}}\) ．
\subsection{题目7}
四．\(y=f(x)\) 为 \([a, b]\) 上严格单调递增的连续函数，\(f(a)=\alpha, f(b)=\beta\) 。证明：\(f(x)\) 的值域为 \([\alpha, \beta]\) ，且 \(x=f^{-1}(y)\) 为连续函数．
\subsection{题目8}
五．曲面 \(\Omega: x^{4}+y^{4}+z^{4}=1, V\) 是 \(\Omega\) 所围区域的体积．求证：\(V=\frac{1}{3} \iint_{\Omega} \frac{\mathrm{d} S}{\sqrt{x^{6}+y^{6}+z^{6}}}\) ．
\subsection{题目9}
六．\(f \in C^{2}[a, b], f_{x}(a)=f_{x}(b)=0\) ，求证： \(\int_{a}^{b}\left|f_{x}\right|^{6} \mathrm{~d} x \leq 25 \int_{a}^{b}\left|f_{x}\right|^{2}\left|f_{x x}\right|^{2} \mathrm{~d} x \cdot\left(\max _{a \leq x \leq b}|f(x)|\right)^{2}\) ．
\end{document}
